\documentclass[tikz,border=6pt]{standalone}
\usepackage{ctex,amsmath}
\usetikzlibrary{arrows.meta,fit,backgrounds}
\usepackage{../../mymath}
\begin{document}
\begin{tikzpicture}[>=Latex,font=\small,
 box/.style={draw=blue!60!black,fill=blue!5,rounded corners,align=center,minimum height=7mm,text width=35mm},
 norm/.style={box,fill=orange!12,minimum height=6mm},
 add/.style={circle,draw,fill=white,inner sep=0pt,minimum size=5mm},
 note/.style={align=center,font=\footnotesize}]
\node[note] (src) at (0,0) {源词元 $x_{1:S}$，$[B,S]$};
\node[note] (tgt) at (7,0) {右移目标 $(\mathrm{BOS},y_1,\ldots,y_{T-1})$\\$[B,T]$};
\foreach \x/\p in {0/e,7/d}{
 \node[box] (\p emb) at (\x,1) {词嵌入 $\times\sqrt d$};
 \node[add] (\p pos) at (\x,1.8) {$+$};
 \node[note] (\p pe) at (\x-2,1.8) {位置编码};
 \draw[->,thick] (\p emb)--(\p pos);
 \draw[->,thick] (\p pe)--(\p pos);
}
\draw[->,thick] (src)--(eemb);\draw[->,thick] (tgt)--(demb);
\node[box] (esa) at (0,3.1) {多头自注意力\\$\mat{Q},\mat{K},\mat{V}\leftarrow \mat{C}_{\ell-1}$};
\node[box] (eff) at (0,6.1) {逐位置前馈网络\\$d\to f\to d$，ReLU};
\node[box] (dsa) at (7,3.1) {掩码多头自注意力\\$\mat{Q},\mat{K},\mat{V}\leftarrow \mat{X}_{\ell-1}$};
\node[box] (dca) at (7,6.1) {多头交叉注意力\\$\mat{Q}\leftarrow \mat{U}_\ell$，$\mat{K},\mat{V}\leftarrow \mat{C}$};
\node[box] (dff) at (7,9.1) {逐位置前馈网络\\$d\to f\to d$，ReLU};
\foreach \p/\x/\y in {esa/0/3.1,eff/0/6.1,dsa/7/3.1,dca/7/6.1,dff/7/9.1}{
 \coordinate (\p in) at (\x,\y-0.7);
 \node[add] (\p add) at (\x,\y+0.8) {$+$};
 \node[norm] (\p ln) at (\x,\y+1.6) {LayerNorm};
 \draw[->,thick] (\p in)--(\p.south);
 \draw[->,thick] (\p.north)--(\p add);
 \draw[->,thick] (\p add)--(\p ln);
 \draw[->,thick,gray!80!black] (\p in)--++(2.2,0)|-(\p add.east);
 \fill (\p in) circle (1.2pt);
}
\draw[->,thick] (epos)--node[left,font=\footnotesize] {$\mat{C}_0$} (esain);
\draw[->,thick] (dpos)--node[right,font=\footnotesize] {$\mat{X}_0$} (dsain);
\draw[->,thick] (esaln)--(effin);
\draw[->,thick] (dsaln)--node[left,font=\footnotesize] {$\mat{U}_\ell$} (dcain);
\draw[->,thick] (dcaln)--node[left,font=\footnotesize] {$Z_\ell$} (dffin);
\node[note] (mem) at (0,8.7) {最终编码器记忆\\$\mat{C}=\mat{C}_{N_e}$，$[B,S,d]$};
\draw[->,thick] (effln)--(mem);
\draw[->,thick,blue!70!black] (mem.east)--(3.6,8.7)--(3.6,6.1)--node[above,font=\footnotesize] {$\mat{K},\mat{V}$} (dca.west);
\node[note,text width=24mm] at (3.6,9.6) {同一份源记忆\\送入每个解码层};
\node[box] (linear) at (7,12) {线性词表投影\\$[B,T,\mat{V}_{\mathrm{vocab}}]$};
\node[box] (softmax) at (7,13.2) {Softmax（词表维）};
\node[note] (out) at (7,14.3) {$p(y_t\mid y_{<t},x_{1:S})$\\训练时与 $y_{1:T}$ 对齐};
\draw[->,thick] (dffln)--node[right,font=\footnotesize] {$\mat{X}_{N_d}$} (linear);
\draw[->,thick] (linear)--(softmax);\draw[->,thick] (softmax)--(out);
\begin{scope}[on background layer]
 \node[draw=gray,dashed,fit=(esa)(esain)(effln),inner xsep=7mm,inner ysep=3mm] (enc) {};
 \node[draw=gray,dashed,fit=(dsa)(dsain)(dffln),inner xsep=7mm,inner ysep=3mm] (dec) {};
\end{scope}
\node[note,anchor=south] at (enc.north) {编码器层 $\times N_e$};
\node[note,anchor=south] at (dec.north) {解码器层 $\times N_d$};
\node[note,text width=44mm,anchor=north] at (0,14) {原始 Transformer：Post-LN\\[4pt]
源侧自注意力：双向可见；\\目标侧自注意力：因果可见；\\交叉注意力：读取有效源位置。\\[4pt]
灰色旁路：残差连接；\\圆圈 $+$：逐元素相加；\\相加后进行 LayerNorm。\\[4pt]
图中省略 Dropout 与填充掩码。\\堆叠层结构相同，参数各自独立。};
\end{tikzpicture}
\end{document}
